% -----------------------------------------------------------
% -----------------------------------------------------------
\part{Book Notes}
% -----------------------------------------------------------
% -----------------------------------------------------------
	\label{BN}
	
% -----------------------------------------------------------
\section{\textit{Slow Productivity} (Cal Newport)}
% -----------------------------------------------------------

	See my one-page summary notes \href{https://drive.google.com/file/d/1QVwVWPOWDxaMUlvx2ybjXujJ8yDLyyFz/view?usp=sharing}{here}. The notes below are just a copy-paste from those.
	
	\begin{enumerate}
	
		\item \textbf{\tr{\textsc{Do fewer things.}}}
			\begin{compactenum}
				\item ``Strive to reduce your obligations to the point where you can easily imagine accomplishing them with time to spare. Leverage this reduced load to more fully embrace and advance the small number of projects that matter most'' (46--47).
				\item Limit the number of overall projects you have. This comes from limiting missions.
				\item \textit{Overhead tax:} The ongoing administrative tasks, like email and meetings, associated with certain projects.
				\item Minimize overhead tax by choosing projects carefully.
				\item Don't accept any ``task engines,'' or projects that generate tasks and overhead taxes endlessly.
				\item Slack is the enemy.
			\end{compactenum}
			
		\item \textbf{\tr{\textsc{Work at a natural pace.}}}
			\begin{compactenum}
				\item ``Don't rush your most important work. Allow it instead to unfold along a sustainable timeline, with variations in intensity, in settings conducive to brilliance'' (94--95).
				\item Fit creative work into a flourishing existence---not the other way around.
				\item Your effort is not an obstacle to a life well-lived---it is part of it.
			\end{compactenum}
			
		\item \textbf{\tr{\textsc{Obsess over quality.}}}
			\begin{compactenum}
				\item ``Obsess over the quality of what you produce, even if this means missing opportunities in the short term. Leverage the value of these results to gain more and more freedom in your efforts over the long term'' (138).
				\item ``Don't ever write for radio.''
				\item ``You should be focused on the quality of what you produce because quality turns out to be connected in unexpected ways to our desire to escape pseudo-productivity and embrace something slower'' (140).
			\end{compactenum}
			
		\item \textbf{\tr{\textsc{The cognitive effort required for significant achievements is \textit{not} busy-ness; in fact, that effort is \textit{incompatible} with busy-ness.}}}
			\begin{compactenum}
				\item ``Obsessing over quality often demands that you slow down,vas the focus required to get better is simply not compatible with busyness'' (140).
				\item ``Working with unceasing intensity is artificial and unsustainable...it estranges us from our fundamental nature, generates misery, and, from a strictly economic perspective, almost certainly holds us back from reaching our full capabilities'' (100).
			\end{compactenum}
		
		\item \textbf{You must take a long-term view of your efforts. The goal is legacy-building accomplishments, which take time to develop.}
			\begin{compactenum}
				\item Real productivity is ultimately measured in the long-term, over long periods of time.
				\item ``Modest results over time. The path is long.''
				\item ``Because if you put a drop in a bucket every day, after three hundred and sixty-five days, the bucket's going to have some water in it...This is what ultimately matters: where you end up, not the speed at which you get there, or the number of people you impress with your jittery busyness along the way'' (172).
			\end{compactenum}
	
		\item \textbf{You must to define measures of actual productivity for any project you work on.}
			\begin{compactenum}
				\item \textit{Pseudo-productivity:} ``The use of visible activity as the primary means of approximating actual productive effort'' (23--24).
				\item Understand what pseudo-productivity is for any given project.
				\item Work doesn't count unless it's actual productivity---deep work on meaningful projects you have chosen.
			\end{compactenum}
			
		\item \textbf{Irresponsibility requires eternal vigilance.}
			\begin{compactenum}
				\item Always quiet-quit. Nothing extra.
			\end{compactenum}
			
		\item \textbf{Planfully approach each project, and don't rush.}
			\begin{compactenum}
				\item You have to make time each day for slow productivity and deep work.
				\item Have a weekly planning session.
				\item If you are planning and trying your best, you are progressing. Be ok with going slow.
			\end{compactenum}
			
		\item \textbf{Develop rituals and locations associated with deep work and the right kind of distraction-free focus.}
			\begin{compactenum}
				\item Find and develop seasonal rituals and practices to allow yourself to work in a natural rhythm.
				\item Rituals and locations are for inducing certain psychological states.
				\item The power of rituals ``is found not in the specifics of their activities but in the transformative effect these activities have on the mind. The more striking and notable the behaviors, the better chance they have of inducing useful changes'' (130).
				\item ``Ensure your rituals are sufficiently striking to effectively shift your mental state into something more supportive of your goals'' (131).
				\item What are the physical spaces where you work? How do they affect you? 
				\item How can you configure those physical spaces to benefit you?
			\end{compactenum}
			
		\item \textbf{You must commit to doing something \textit{very well}; you must obsess over quality.}
			\begin{compactenum}
				\item ``All of us who do creative work, we get into it because we have good taste...But it's like there's a gap. That for the first couple years that you're making stuff, what you're making, isn't so good...it's not quite that good...If you're just starting off and you're entering into that phase, you gotta know it's totally normal and the most important possible thing you can do is do a lot of work...Put yourself on a deadline so that every week or every month you know you're gonna finish one story...It's only by actually going through a volume of work that you're actually going to catch up and close that gap, and the work you're making will be as good as your ambitions'' (146).
				\item ``To those just beginning their careers...keep putting in the work, as it's only through this deliberate effort that the gap will close'' (147).
				\item ``When you spend two years reading and critiquing and admiring work by other young writers pushing their prose in new and interesting directions, your standards for what writing can achieve sharpen....there's a reason why MFA programs are so common among successful writers: they provide an effective training regimen for literary taste'' (147--148).
			\end{compactenum}
			
		\item \textbf{Develop your taste by intaking high-quality material from many different fields as often as you're able.}
			\begin{compactenum}
				\item Getting to know fields other than your own is crucial for creativity.
				\item Learn to develop good taste in other fields.
			\end{compactenum}
			
		\item \textbf{In all your projects, bet on yourself---this means, be confident enough to push them forward boldly.}
			\begin{compactenum}
				\item ``Betting on yourself in this manner---with nontrivial stakes for failure but attractive rewards for success---is a good general strategy for pushing the quality of your work to a new level''(160).
				\item You are either capable, or can become capable.
			\end{compactenum}
			
		\item \textbf{In what ways are digital devices stopping or distracting you?}
			\begin{compactenum}
				\item Devices are instruments of pseudo-productivity and overhead taxes.
				\item Configur your computer to minimize distractions, including email and text notifications.
			\end{compactenum}
	
	\end{enumerate}
	
% -----------------------------------------------------------
\pagebreak{}
% -----------------------------------------------------------